%===============================================================================
% SEÇÃO 3 — Revisão da Literatura de Base
%===============================================================================

\section{Revisão da Literatura de Base}
\label{sec:revisao}

\subsection{Controle de processos endotérmicos}

Reatores tubulares aquecidos externamente, como os utilizados em processos de
reforma a seco, apresentam forte acoplamento entre o balanço de energia no
forno e o perfil de temperatura ao longo do leito catalítico, resultando em
dinâmica lenta e não linear.\citep{luyben2014} Uma abordagem comum em
aplicções industriais é o uso de estruturas de controle em cascata, nas quais
uma malha interna regula a variável associada ao fornecimento de calor enquanto
uma malha externa atua sobre a temperatura de saída do reator ou sobre um ponto
representativo do leito.\citep{luyben2014} Essa configuração melhora a
rejeição a perturbações de carga e de composição e reduz o efeito dos atrasos
associados à inércia térmica do conjunto forno-reator.

Estudos recentes em reforma a seco de metano com CO\(_2\) apontam que variações
na razão de alimentação e nos perfis de aquecimento podem alterar
consideravelmente a sensibilidade do processo, motivando o uso de técnicas
avançadas como estruturas neuro-fuzzy (ANFIS) e controle preditivo baseado em
modelos (MPC).\citep{essien2019} Em qualquer dessas abordagens, a qualidade
das medições de temperatura e a localização adequada dos sensores ao longo do
reator são decisivas para o sucesso da estratégia de controle.

\subsection{Análise termodinâmica como suporte ao controle}

Análises termodinâmicas de equilíbrio têm sido amplamente empregadas para
mapear as janelas operacionais da DRM e da BDR em função de temperatura,
presão e razão CO\(_2\)/CH\(_4\).\citep{thermodynamics2025,lavoie2014} Esses
estudos identificam as regiões em que a conversão de CH\(_4\) e CO\(_2\) é
elevada e a formação de carbono sólido é termodinamicamente desfavorecida,
fornecendo limites superiores e inferiores diretos para as variáveis de
operação. Essas informações são valiosas tanto no projeto de malhas
regulatórias quanto no desenvolvimento de esquemas de proteção e, em níveis
hierárquicos mais elevados, em problemas de otimização em tempo real e
controle preditivo.\citep{thermodynamics2025}
